% =============================================================================
% ch04_theory_impossibility.tex — The Unknowable Regime
% Long-form expansion of kbound.tex §5.1 (Result 1: non-identifiability,
% witness, Le Cam two-point bound) and §5.2 (Result 2: Bayes gate, plug-in
% regret, minimax floor). Theorem/proposition statements are VERBATIM from
% the paper; proofs are expanded line by line; numerical-validation numbers
% are quoted from experiments/kbound/theory_validation/. NO preamble here.
% =============================================================================

\chapter{The Unknowable Regime: Non-Identifiability and the Minimax Floor}\label{ch:theory1}

The first half of the theory draws the boundary of the possible. It does so on
ground the literature already owns: that label-free identification of a
shift-dependent quantity is impossible without an assumption on the shift is due to
Garg et al.\ \cite{garg2022atc}, and the decision theory of abstaining versus
predicting under shift --- with an observable-quantity-versus-shift-budget threshold
--- to Kalai and Kanade \cite{kalai2021abstain}. This chapter does not re-derive
that general impossibility; it \emph{specializes and sharpens} it to the sign of the
adaptation benefit, and proves two results in that specialized form. The first
(Theorem~\ref{thm:imp}, with its quantitative Le~Cam strengthening,
Theorem~\ref{thm:t1-impquant}) shows that the unknowable regime of
Definition~\ref{def:form-regimes} is not an artifact of weak evidence
features or small batches: there exist pairs of target worlds whose
label-free evidence laws coincide \emph{exactly} --- for every evidence map
and every unlabeled sample size --- while their adaptation benefits have
opposite signs, and on such pairs the optimal committal error admits a
closed-form floor, $1-\TV$, which equals $1$ at the witness. That closed form is our
sharpening of the general hardness to an exact identity for this question. The
second (Theorem~\ref{thm:gate}) prices every violation of the boundary: the
regret-to-oracle of any plug-in gate equals, \emph{exactly}, the expected
benefit forfeited on the evidence values where the estimated sign is wrong,
which yields the matching minimax floor $\E[|\Delta(Z)|]/2$ for the
unknowable regime and justifies abstention on the low-margin band.

The paper's running example sets the stakes. Two hospitals' scanners induce
identical confidence and entropy statistics in a deployed classifier, yet in
one hospital adapting helps and in the other --- a flipped label frequency
--- it badly hurts; no label-free rule can separate them, so the safe action
is to abstain. The mathematics below shows that this story is not a corner
case but a face of the problem: whenever the posterior $P_T(Y\mid X)$ can
move while the covariate law $P_T(X)$ does not, the evidence
$\sigma$-algebra simply does not contain the answer. Everything in this
chapter is proved against the strongest reading of the protocol of
Chapter~\ref{ch:formulation}: arbitrary evidence maps, arbitrary
(randomized) rules, unlimited unlabeled data, and free access to labeled
source data.

We emphasize the chapter's intellectual posture at the outset, because the
paper does. The general impossibility is not ours: label-free identification of
target performance is as hard as identifying the optimal predictor and needs an
assumption on the shift \cite{garg2022atc}, and the abstain-versus-predict decision
theory under shift is already established \cite{kalai2021abstain}; more broadly,
non-identifiability of unlabeled target \emph{risk} is close to folklore --- it is
the reason the unsupervised risk estimation, domain adaptation, and label-shift
literatures all carry structural assumptions \cite{bendavid2010theory,
mansour2009domain, lipton2018detecting}. What this chapter contributes is not the
impossibility but its \emph{specialization} to the adapt/freeze decision: (i) the
statement at the level of the \emph{decision}, where it justifies a specific third
action --- the benefit-sign instance of the under-shift abstention rule of
\cite{kalai2021abstain}; (ii) an explicit, maximally separated witness pair on which
every later experiment can be run; (iii) the exact minimax form of the floor via
the Le~Cam two-point method \cite{lecam1973convergence,
tsybakov2009introduction}, sharpening the general hardness to an identity for this
question; and (iv) the exact regret identity that converts the floor into a
deployment-relevant quantity. Each component is validated numerically by a
machine-checked script in the public repository \cite{kbound2026repo}, and we quote
those checks where they occur.

\section{The non-identifiability theorem}\label{sec:t1-imp}

We begin with the qualitative statement, copied verbatim from the paper.

\begin{theorem}[Non-identifiability; with witness]\label{thm:imp}
Let $g(Z)\in\{\textsc{adapt},\textsc{freeze},\textsc{abstain}\}$ be any
(possibly randomized) rule depending on label-free evidence $Z$. There exist
two target worlds $P_T^1,P_T^2$ with
\textup{(i)} $\Law(Z\mid P_T^1)=\Law(Z\mid P_T^2)$ and
\textup{(ii)} $\sign\Delta^1=-\sign\Delta^2\neq 0$.
For any such pair, $\Law(g)$ is identical across the two worlds; hence no
label-free rule commits to the benefit-maximizing action in both, and the
action minimizing the worst-case committal regret over $\{P_T^1,P_T^2\}$ is
\textsc{abstain}.
\end{theorem}

Before the proof, three reading notes. First, the theorem has an
existence half and a universal half: it \emph{exhibits} one pair of worlds
satisfying (i)--(ii) (the witness), and then shows that for \emph{every}
pair satisfying (i)--(ii) --- witness or otherwise --- every rule is stuck.
Second, the conclusion is about laws, not samples: $\Law(g)$ identical
across worlds means the rule's behavior is indistinguishable \emph{in
distribution}, so no accounting trick (more batches, more features, more
randomization) recovers a difference. Third, the optimality of
\textsc{abstain} is a genuine minimax statement over a two-point family,
in the classical mold of two-point lower bounds
\cite{lecam1973convergence, tsybakov2009introduction}; the next section
upgrades it to an exact quantitative floor.

\begin{proof}
The proof has four steps: construct the witness; verify properties (i) and
(ii); show every label-free rule has world-independent law; and solve the
resulting worst-case decision problem.

\emph{Step 1 (witness construction; non-vacuous).}
Let $X\sim\mathcal N(0,1)$ under \emph{both} worlds, so $P_T^1(X)=P_T^2(X)$.
Take the fixed predictors $f_0(x)=\mathbf 1[x>0]$ and
$f_a(x)=\mathbf 1[x<0]$, and $0/1$ loss. Define the two conditional label
laws degenerately:
\[
P_T^1(Y\mid X)\colon\ Y=\mathbf 1[X>0],
\qquad
P_T^2(Y\mid X)\colon\ Y=\mathbf 1[X<0].
\]
In world~1, $f_0(X)=\mathbf1[X>0]=Y$ holds almost surely, so
$R_T^1(f_0)=\Pr[f_0(X)\neq Y]=0$; and since $\Pr[X=0]=0$, we have
$f_a(X)=\mathbf1[X<0]=1-Y$ almost surely, so $R_T^1(f_a)=1$. Hence
\[
\Delta^1=R_T^1(f_0)-R_T^1(f_a)=0-1=-1<0,
\]
and freezing is the benefit-maximizing action in world~1. World~2 is the
mirror image: $f_a(X)=Y$ and $f_0(X)=1-Y$ almost surely, so
$R_T^2(f_0)=1$, $R_T^2(f_a)=0$, and $\Delta^2=+1>0$: adapting is correct.
The benefit gap is maximal ($|\Delta^j|=1$, the full range of $0/1$ risk),
so the construction is not a degenerate tie.

\emph{Step 2 (identical evidence laws).}
The evidence $Z=\phi(X_{1:m},f_0,f_a)$ is, for any evidence map $\phi$ and
any batch size $m$, a function of $X_{1:m}$ and of the fixed maps $f_0,f_a$
only: writing $\psi(\cdot):=\phi(\cdot,f_0,f_a)$, we have
$Z=\psi(X_{1:m})$ with $\psi$ measurable and world-independent. The law of
$X_{1:m}$ is $\mathcal N(0,1)^{\otimes m}$ in both worlds, so the
pushforward laws coincide:
$\Law(Z\mid P_T^1)=\Law(X_{1:m})\circ\psi^{-1}=\Law(Z\mid P_T^2)$. This is
property (i); property (ii) was verified in Step~1
($\sign\Delta^1=-1=-\sign\Delta^2\neq0$). Thus (i) and (ii) hold
simultaneously.

\emph{Step 3 (the rule's law is world-independent).}
A label-free rule is a measurable map of the evidence and of independent
randomness: $g=\Gamma(Z,U)$, where $U$ is a randomization variable whose
law does not depend on the target world (it is generated by the system, not
by the data), and $\Gamma$ is measurable. Under world $j$ the law of $g$ is
the pushforward of $\Law(Z\mid P_T^j)\otimes\Law(U)$ under $\Gamma$. By
(i) the two product laws coincide, hence so do the laws of $g$. In
particular the three numbers
\[
q:=\Pr[g=\textsc{adapt}],\qquad r:=\Pr[g=\textsc{freeze}],\qquad
1-q-r=\Pr[g=\textsc{abstain}]
\]
are the \emph{same} in both worlds.

\emph{Step 4 (no simultaneous correctness, and the worst-case-optimal
action).}
The benefit-maximizing action is \textsc{freeze} in world~1 and
\textsc{adapt} in world~2. A rule committing to the correct action with
probability one in world~1 needs $r=1$; in world~2 it needs $q=1$; both
cannot hold since $q+r\le1$. Hence no label-free rule commits to the
benefit-maximizing action in both worlds. To identify the minimax action,
measure committal regret as the excess $0/1$ risk of a committed action
versus the better committed action in that world, and assign abstention a
separate coverage cost (it is a different kind of outcome --- a flagged
non-decision --- and is priced separately; the quantitative theorem below
makes this explicit). In world~1, committing to \textsc{adapt} costs
$R_T^1(f_a)-R_T^1(f_0)=|\Delta^1|=1$ and committing to \textsc{freeze}
costs $0$; so world~1 charges expected committal regret
$q\cdot|\Delta^1|=q$. Symmetrically world~2 charges
$r\cdot|\Delta^2|=r$. The worst case over the two-point family is
\[
W(q,r)\;=\;\max(q,\,r),
\]
to be minimized over the simplex $\{q,r\ge0,\ q+r\le1\}$. Since
$W(q,r)\ge0$ with equality if and only if $q=r=0$, the unique minimizer is
$q=r=0$, i.e.\ the rule that plays \textsc{abstain} with probability one.
\end{proof}

The paper attaches an explicit calibration of the theorem's weight, which
we reproduce because it governs how the result should be cited.

\begin{remark}[Weight of the result]\label{rem:t1-weight}
Theorem~\ref{thm:imp} is close to known non-identifiability of target risk
without labels; its purpose here is to \emph{justify the abstain action},
not to claim impossibility as novel. The quantitative strengthening below
(Theorem~\ref{thm:t1-impquant}) makes it a tight minimax statement.
\end{remark}

\begin{remark}[Both hypotheses are necessary]\label{rem:t1-hypotheses}
The two-point structure is sharp in both coordinates. If (i) fails --- the
evidence laws differ by total variation $\tau>0$ --- then
Theorem~\ref{thm:t1-impquant} shows the optimal committal error drops to
$1-\tau<1$, and a certificate may commit on the distinguishable part
(Chapter~\ref{ch:theory2} builds exactly that). If (ii) fails --- the two
worlds' benefits share a sign --- then committing to the common
benefit-maximizing action is correct in both worlds even though the
evidence cannot tell them apart; indistinguishability is harmless when the
decision does not depend on the distinction. The impossibility requires the
conjunction: \emph{the evidence cannot separate the worlds, and the worlds
disagree about what to do.} The unknowable regime of
Definition~\ref{def:form-regimes} is precisely the set of evidence values
where some such pair lives.
\end{remark}

\section{The explicit witness, in full}\label{sec:t1-witness}

Because the witness carries both the theorem's existence half and the
paper's cleanest experiment, we devote a section to it: first the algebra
in complete detail, then its interpretation, then the empirical
realization.

\subsection{Construction and verification}\label{sec:t1-witness-algebra}

The witness consists of: covariate law $X\sim\mathcal N(0,1)$, shared by
both worlds; predictors $f_0(x)=\mathbf1[x>0]$ and $f_a(x)=\mathbf1[x<0]$,
fixed and known; loss $0/1$; and the two degenerate posteriors
$Y=\mathbf1[X>0]$ (world~1) and $Y=\mathbf1[X<0]$ (world~2). Three features
deserve explicit verification beyond the proof above.

\emph{The risks, entry by entry.} Under world~1, the event
$\{f_0(X)\neq Y\}=\{\mathbf1[X>0]\neq\mathbf1[X>0]\}=\emptyset$, so
$R_T^1(f_0)=0$ exactly (not merely almost). The event
$\{f_a(X)\neq Y\}=\{\mathbf1[X<0]\neq\mathbf1[X>0]\}=\{X\neq0\}$ has
probability one under the atomless Gaussian, so $R_T^1(f_a)=1$. The four
risks are
\[
R_T^1(f_0)=0,\quad R_T^1(f_a)=1,\quad R_T^2(f_0)=1,\quad R_T^2(f_a)=0,
\]
giving $\Delta^1=-1$ and $\Delta^2=+1$: the largest benefit gap the $0/1$
scale permits, in opposite directions.

\emph{Evidence-law equality, for every $\phi$ and every $m$.} The claim is
not that some chosen panel of statistics happens to match across worlds; it
is that the entire observable $\sigma$-algebra matches. Any label-free
evidence is $Z=\psi(X_{1:m})$ for a fixed measurable $\psi$ (the models
being fixed known maps, they are absorbed into $\psi$), and equality of the
laws of $X_{1:m}$ forces equality of the laws of $Z$ --- a pushforward of
equal measures under one map. Consequently \emph{no} statistic computable
without target labels --- no entropy histogram, no drift score, no
disagreement pattern, no density-ratio diagnostic, nor any statistic anyone
will ever invent --- distinguishes the worlds at any batch size. The
quantifier order is the point: the construction defeats all evidence maps
simultaneously, not a fixed one.

\emph{What kind of shift this is.} Between the two worlds the covariate law
is untouched and the posterior flips: $P_T^1(Y{=}1\mid x)=\mathbf1[x>0]$
versus $P_T^2(Y{=}1\mid x)=\mathbf1[x<0]$. In the taxonomy of shifts this
is pure concept shift \cite{quinonero2009dataset}, the row of the
knowability hierarchy that Chapter~\ref{ch:theory2} marks non-identifiable;
Example~\ref{ex:form-flip} of Chapter~\ref{ch:formulation} is its smoothed
sibling, with overlapping classes and $|\Delta|=2\Phi(1)-1\approx0.683$ in
place of the extremal $|\Delta|=1$, showing the phenomenon survives label
noise. It is also instructive to view the witness through the
disagreement-region lens of Theorem~\ref{thm:disagree}: here the observable
disagreement region is $D=\{x:f_0(x)\neq f_a(x)\}=\{x\neq0\}$, essentially
the whole line, and the candidate's accuracy on it is $a_a^D=0$ in world~1
and $a_a^D=1$ in world~2 --- the identity
$\sign\Delta=\sign(a_a^D-\tfrac12)$ holds in both worlds, and what the
witness makes impossible is precisely any label-free \emph{bracketing} of
$a_a^D$ around $\tfrac12$. The impossibility and the
sign-of-difference characterization are two views of one boundary.

\subsection{Empirical realization}\label{sec:t1-witness-empirical}

The witness is constructive, so it can be run. The paper reports a clean
witness experiment (the full experimental protocol appears in
Chapter~\ref{ch:exp1}): two simulated worlds with $X\sim\mathcal N(0,1)$,
$f_0=\mathbf1[X{>}0]$, $f_a=\mathbf1[X{<}0]$, and flipped labels, so the
benefit is exactly opposite (world means $-1.0$ versus $+1.0$) while every
evidence feature is a function of $X$ only. Across $300$ instances, all
five evidence features used by \KGA\ are statistically indistinguishable
between the worlds --- two-sample Kolmogorov--Smirnov $p$-values in
$[0.16,0.96]$, all above $0.05$ --- yet the ground truth is opposite, so
$Z$ cannot decide. \KGA\ abstains on $100\%$ of instances, and a rule
forced to commit by $\sign\widehat\Delta$ pays regret $0.51$, near chance.
This is the clean theory--experiment bridge
(Figure~\ref{fig:t1-witness}).

The repository's theorem validator
(\texttt{experiments/kbound/theory\_validation/val\_thm1\_lecam.py}, results
in \texttt{results\_thm1\_lecam.json} \cite{kbound2026repo}) re-runs the
witness at scale with $20{,}000$ Monte-Carlo batches of $n=64$ unlabeled
samples per world and reports four independent checks. (1)~The closed-form
benefits $\Delta^1=-1$, $\Delta^2=+1$ are reproduced exactly by Monte
Carlo. (2)~For each of five per-sample evidence features (the raw
covariate, the margin $|x|$, the predicted class of $f_0$, the
disagreement indicator, and a smooth confidence transform), the
permutation-debiased total-variation estimate between the two worlds'
evidence laws is at most $1.9\times10^{-3}$ --- zero up to the estimator's
own null floor --- with two-sample KS $p$-values between $0.28$ and $1.0$;
and the $n$-sample aggregated evidence statistic has exact TV $0$ for every
$n$ in $\{1,2,4,\dots,256\}$ (debiased estimates $\le4.8\times10^{-3}$).
(3)~Optimizing the committal error over a dense family of threshold rules
(all thresholds, both polarities) on each evidence statistic yields
$\inf_g M(g)$ between $0.991$ and $1.000$ across the five statistics,
against the theoretical floor of $1$. (4)~The actual \KGA\
conformal three-way rule, run exactly as in the paper with $\alpha=0.1$,
abstains on $4{,}000$ of $4{,}000$ instances ($100\%$ abstention; zero
false adapts and zero false freezes): its calibrated radius
($\varepsilon\approx1.15$) straddles zero because the evidence carries no
benefit signal, so the certificate cannot fire --- the behavior that
Chapter~\ref{ch:theory2} proves is mandatory for any rule with valid error
control. The paper's summary sentence is exact: in the witness the
empirical $\inf_g M$ is ${\approx}1$ (\KGA\ abstains $100\%$); in the
detectable regime the empirical minimax error tracks the closed-form
$1-\TV$ across all $n$ and shift magnitudes.

\begin{figure}[t]\centering
\includegraphics[width=0.72\textwidth]{fig_witness_clean.png}
\caption[Clean non-identifiability witness]{The clean non-identifiability
witness of Theorem~\ref{thm:imp}, realized empirically: identical evidence
law across the two worlds (two-sample KS $p>0.05$ on all five evidence
features), exactly opposite true benefit ($-1.0$ versus $+1.0$), and
$100\%$ abstention by the \KGA\ rule; a rule forced to commit pays regret
$0.51$, near chance. The experiment instantiates hypotheses (i) and (ii)
of the theorem exactly rather than approximately, which is what makes it a
\emph{witness} rather than a hard benchmark.}
\label{fig:t1-witness}
\end{figure}

\section{The quantitative floor: a Le Cam two-point bound}\label{sec:t1-lecam}

Theorem~\ref{thm:imp} says committal rules fail on matched evidence; it
does not yet say how badly, nor what happens as the evidence laws begin to
separate. The paper's quantitative strengthening answers both with an exact
identity, in the two-point tradition of Le~Cam \cite{lecam1973convergence};
the intuition, in the paper's words: the more alike two such worlds look in
the observable evidence (smaller TV distance), the closer any rule forced
to commit is pushed toward a coin flip, and the error floor vanishes only
as their evidence separates.

We first isolate the two classical ingredients as lemmas, with proofs,
since the theorem's proof is a direct composition of them.

\begin{lemma}[Testing affinity identity]\label{lem:t1-affinity}
Let $Q_1,Q_2$ be probability measures on a common measurable space. Over
all randomized tests $\psi$ (measurable functions with values in $[0,1]$),
\[
\inf_{\psi}\;\bigl\{\E_{Q_1}[\psi]+\E_{Q_2}[1-\psi]\bigr\}
\;=\;1-\TV(Q_1,Q_2),
\]
and the infimum is attained by the deterministic likelihood-ratio test
$\psi^\star=\mathbf1[\mathrm dQ_2>\mathrm dQ_1]$.
\end{lemma}

\begin{proof}
Let $\nu=Q_1+Q_2$, which dominates both measures, and let $q_1,q_2$ be the
corresponding densities. For any test $\psi$,
\[
\E_{Q_1}[\psi]+\E_{Q_2}[1-\psi]
=\int\psi\,q_1\,\mathrm d\nu+\int(1-\psi)\,q_2\,\mathrm d\nu
=1+\int\psi\,(q_1-q_2)\,\mathrm d\nu,
\]
using $\int q_2\,\mathrm d\nu=1$. The integrand is minimized pointwise over
$\psi(z)\in[0,1]$ by setting $\psi=1$ where $q_1-q_2<0$, $\psi=0$ where
$q_1-q_2>0$, and arbitrarily on the crossing set $\{q_1=q_2\}$ (which
contributes zero); this is exactly $\psi^\star=\mathbf1[q_2>q_1]$, a
deterministic test, so randomization buys nothing. The attained value is
\[
1-\int (q_2-q_1)_+\,\mathrm d\nu .
\]
It remains to identify the integral with total variation. For any
measurable $A$,
$Q_2(A)-Q_1(A)=\int_A(q_2-q_1)\,\mathrm d\nu\le\int(q_2-q_1)_+\,\mathrm
d\nu$, with equality at $A^\star=\{q_2>q_1\}$; hence
$\TV(Q_1,Q_2):=\sup_A|Q_1(A)-Q_2(A)|=\int(q_2-q_1)_+\,\mathrm d\nu$, which
also equals $\tfrac12\int|q_1-q_2|\,\mathrm d\nu$ because
$\int(q_2-q_1)\,\mathrm d\nu=0$ forces the positive and negative parts to
have equal mass. Substituting gives the identity
\cite{lecam1973convergence, tsybakov2009introduction, lehmann2005testing}.
\end{proof}

\begin{lemma}[Data-processing inequality for total variation]\label{lem:t1-dpi}
Let $P_1,P_2$ be probability measures and $T$ a measurable map. Then
$\TV\bigl(P_1\circ T^{-1},\,P_2\circ T^{-1}\bigr)\le\TV(P_1,P_2)$, with
equality whenever $T$ is a sufficient statistic for the pair
$\{P_1,P_2\}$.
\end{lemma}

\begin{proof}
For the inequality, every event of the pushforward laws is the preimage of
an event: $\sup_A|P_1(T^{-1}A)-P_2(T^{-1}A)|$ is a supremum over the
sub-collection $\{T^{-1}A\}$ of measurable sets, hence at most the
supremum over all measurable sets, which is $\TV(P_1,P_2)$. For the
equality claim: if $T$ is sufficient for $\{P_1,P_2\}$, the likelihood
ratio $\mathrm dP_2/\mathrm dP_1$ is ($P_1$-a.s.) a function of $T$ by the
factorization theorem \cite{lehmann2005testing}, so the optimal set
$A^\star=\{\mathrm dP_2>\mathrm dP_1\}$ realizing the supremum in
Lemma~\ref{lem:t1-affinity}'s identification of $\TV$ is itself of the
form $T^{-1}(B)$, and the restricted supremum already attains the
unrestricted one.
\end{proof}

With the two lemmas in hand we state the paper's theorem verbatim.

\begin{theorem}[Quantitative non-identifiability; Le Cam two-point lower bound]
\label{thm:t1-impquant}\label{thm:imp-quant}
Consider two target worlds $P_T^1,P_T^2$ with $\Delta^1<0<\Delta^2$. Let the
label-free evidence be $Z=\phi(X_{1:n},f_0,f_a)$ with $n$-sample evidence
laws $P_{T,Z}^{j,(n)}:=\Law(Z\mid P_T^j)$. For any (randomized) committal
rule $g(Z)\in\{\textsc{adapt},\textsc{freeze}\}$ define the mixed committal
error $M(g):=\Pr_{P_T^1}[g=\textsc{adapt}]+\Pr_{P_T^2}[g=\textsc{freeze}]$.
Then
\[
\inf_{g} M(g)\;=\;1-\TV\!\big(P_{T,Z}^{1,(n)},P_{T,Z}^{2,(n)}\big)\;\ge\;1-\TV\!\big(P_T^{1,(n)},P_T^{2,(n)}\big).
\]
In the witness above ($Z$ a function of $X$, $P_T^1(X)=P_T^2(X)$) the
evidence TV is $0$, so $\inf_g M(g)=1$ for \emph{every} $n$: no committal
rule beats blind guessing in the worst world, and under any positive
coverage cost \textsc{abstain} is uniquely worst-case optimal. If the
worlds become $\delta$-detectable so the evidence laws separate,
$\inf_g M(g)=1-\TV\to0$ at the separation rate (for a location shift
summarized by a sufficient statistic, $\TV=2\Phi(\tfrac12\sqrt n\,\delta)-1$).
\end{theorem}

\begin{proof}
We expand the paper's proof into five steps.

\emph{Step 1 (reduction to binary testing).}
A randomized committal rule on $Z$ is characterized by its adapt-probability
$\psi(z):=\Pr[g=\textsc{adapt}\mid Z=z]\in[0,1]$, a randomized test of the
two evidence laws. Identify $\{g=\textsc{adapt}\}$ with deciding world~2
(where adapting is correct, since $\Delta^2>0$). Writing
$Q_j:=P_{T,Z}^{j,(n)}$, the mixed committal error is
\[
M(g)=\Pr_{Q_1}[g=\textsc{adapt}]+\Pr_{Q_2}[g=\textsc{freeze}]
=\E_{Q_1}[\psi(Z)]+\E_{Q_2}[1-\psi(Z)],
\]
the sum of the test's type-I and type-II errors. The infimum of $M$ over
committal rules is therefore the infimum of the testing risk over
randomized tests.

\emph{Step 2 (the equality).}
Lemma~\ref{lem:t1-affinity} applied to $Q_1,Q_2$ gives
$\inf_g M(g)=1-\TV(Q_1,Q_2)$, attained by the likelihood-ratio rule that
plays \textsc{adapt} exactly on $\{\mathrm dQ_2>\mathrm dQ_1\}$. Note this
is an identity, not merely a bound: the optimal committal rule's mixed
error \emph{equals} one minus the evidence separation. This is the
tightness direction --- no slack is lost between the decision problem and
the total-variation geometry, because the optimal test is exhibited.

\emph{Step 3 (the inequality to the full data).}
The evidence is a measurable reduction of the raw unlabeled batch:
$Z=T(X_{1:n})$ with $T:=\phi(\cdot,f_0,f_a)$. By Lemma~\ref{lem:t1-dpi},
$\TV(Q_1,Q_2)\le\TV\bigl(P_T^{1,(n)},P_T^{2,(n)}\bigr)$, the total
variation between the worlds' $n$-sample \emph{covariate-data} laws.
Negating, $1-\TV(Q_1,Q_2)\ge1-\TV(P_T^{1,(n)},P_T^{2,(n)})$: summarizing
data can only raise the error floor. Equality holds when $Z$ is sufficient
for the two-world testing problem --- the evidence then loses nothing.

\emph{Step 4 (the witness: floor $1$ at every $n$).}
In the witness, $Z$ depends on the worlds only through $\Law(X_{1:n})$,
and these laws coincide; hence $Q_1=Q_2$, $\TV(Q_1,Q_2)=0$, and
$\inf_g M(g)=1$ for every $n$. Unpacking what $M(g)=1$ forces: since
$M(g)=\Pr_1[\text{err}]+\Pr_2[\text{err}]\ge1$ for every committal $g$,
at least one world sees error probability $\ge\tfrac12$ --- no committal
rule beats blind guessing in its worst world, and this does not improve
with $n$. Now add a coverage cost $c\in(0,\tfrac12)$ for abstention, and
recall from the witness that a wrong commitment costs $|\Delta^j|=1$.
A rule with commitment probabilities $(q,r)$ (Step~3 of
Theorem~\ref{thm:imp}'s proof; world-independent by matched evidence) has
worst-case expected cost
\[
\max\bigl(q+ (1-q-r)c,\; r+(1-q-r)c\bigr)
\;=\;\max(q,r)+(1-q-r)\,c
\;\ge\;\tfrac{q+r}{2}+(1-q-r)\,c,
\]
which, as a function of the total commitment mass $s=q+r\in[0,1]$, is
$\ge s/2+(1-s)c$, strictly increasing in $s$ because $c<\tfrac12$. The
minimum is at $s=0$, value $c$: always-\textsc{abstain} is the unique
worst-case-optimal rule.

\emph{Step 5 (the detectable rate).}
Suppose the worlds are $\delta$-detectable through a real-valued
per-sample summary with unit noise scale: the summary is Gaussian with
means separated by $\delta$ and variance $1$ (the canonical location
family; in the validator below, the worlds are
$P^1(X)=\mathcal N(-\delta/2,1)$ and $P^2(X)=\mathcal N(+\delta/2,1)$).
The sample mean $\bar X_n$ is sufficient for the location parameter, so by
Lemma~\ref{lem:t1-dpi} the evidence loses nothing and
$\TV(Q_1,Q_2)=\TV\bigl(\mathcal N(-\tfrac\delta2,\tfrac1n),
\mathcal N(+\tfrac\delta2,\tfrac1n)\bigr)$. For two equal-variance
Gaussians $\mathcal N(\theta_1,\sigma^2)$, $\mathcal N(\theta_2,\sigma^2)$
the densities cross exactly at the midpoint $(\theta_1+\theta_2)/2$, so
the positive part integrates to
$\TV=\Phi\bigl(\tfrac{|\theta_1-\theta_2|}{2\sigma}\bigr)-
\Phi\bigl(-\tfrac{|\theta_1-\theta_2|}{2\sigma}\bigr)
=2\Phi\bigl(\tfrac{|\theta_1-\theta_2|}{2\sigma}\bigr)-1$. With
$|\theta_1-\theta_2|=\delta$ and $\sigma=1/\sqrt n$ this is
$\TV=2\Phi(\tfrac12\sqrt n\,\delta)-1$, whence
$\inf_g M(g)=1-\TV=2\Phi(-\tfrac12\sqrt n\,\delta)\to0$ as $\sqrt
n\,\delta\to\infty$: the committal-error floor decays at the Gaussian
separation rate, and the unknowable regime shrinks exactly as fast as the
evidence separates.
\end{proof}

\begin{remark}[Where the bound is and is not tight]\label{rem:t1-tightness}
Two distinct tightness claims should be kept apart. The first equality,
$\inf_g M=1-\TV$ \emph{on the evidence laws}, is always exact
(Lemma~\ref{lem:t1-affinity} exhibits the optimal test); there is no
asymptotic or constant slack anywhere in the chain from decision problem
to total-variation geometry. The second relation, the inequality to the
\emph{full-data} TV, is an inequality precisely because a poorly chosen
evidence map can discard separation that the raw batch contains; it
collapses to equality when $Z$ is sufficient for the two-world problem
(Lemma~\ref{lem:t1-dpi}), as in the location family of Step~5. The design
lesson for evidence panels is contained in that gap: the floor an
operator actually faces is $1-\TV(\text{evidence})$, which the operator
partially controls through $\phi$, bounded below by
$1-\TV(\text{full data})$, which nature controls. In the witness nature
sets the latter to $1$ and the choice of $\phi$ is irrelevant.
\end{remark}

\subsection{Numerical validation of the floor}\label{sec:t1-lecam-val}

The validator \texttt{val\_thm1\_lecam.py} \cite{kbound2026repo} checks the
identity $\inf_g M(g)=1-\TV$ in both regimes, using $20{,}000$ Monte-Carlo
evidence instances per world ($n=64$ unlabeled samples each, seed
$12345$) and optimizing committal rules over a dense
threshold-and-polarity family on each evidence statistic.

In the \emph{witness} regime the results were quoted in
Section~\ref{sec:t1-witness-empirical}: empirical floors
$0.991$--$1.000$ against the theoretical $1$, exact $n$-sample evidence
TV equal to $0$ at every $n$, and $100\%$ \KGA\ abstention. In the
\emph{detectable} regime, $P^1(X)=\mathcal N(-\mu,1)$ and
$P^2(X)=\mathcal N(+\mu,1)$ with opposite benefits, the closed form of
Step~5 (with $\delta=2\mu$) is traced across both the shift magnitude and
the sample size. Sweeping $\mu$ at $n=64$: at $\mu=0.1$ the exact floor is
$1-\TV=0.4237$ and the best empirical committal error is $0.4189$; at
$\mu=0.25$, floor $0.0455$ versus empirical $0.0451$; at $\mu=0.5$ the
floor is $6.3\times10^{-5}$ and the empirical error is $0.0$. Sweeping $n$
at fixed $\mu=0.25$, the (floor, empirical) pairs are
$(0.8026,\,0.8134)$ at $n=1$, $(0.6171,\,0.6055)$ at $n=4$,
$(0.3173,\,0.3209)$ at $n=16$, $(0.0455,\,0.0452)$ at $n=64$, and
$(6.3\times10^{-5},\,0.0)$ at $n=256$ --- the empirical optimum tracks the
closed form throughout, fluctuating around it in both directions by
Monte-Carlo error (the threshold is optimized on the same finite sample,
so small dips below the population floor are expected and honest).
Meanwhile the \KGA\ three-way rule abstains on $100\%$ of instances at
$\mu\in\{0,0.1\}$, on $0.45\%$ at $\mu=0.25$, and on $0\%$ at
$\mu\ge0.5$ with zero committal error: the rule stops abstaining exactly
as the worlds become detectable, which is the operational content of the
phase diagram's $x$-axis (Figure~\ref{fig:form-phase}).

\begin{figure}[t]\centering
\includegraphics[width=0.66\textwidth]{fig_lecam_bound.png}
\caption[Le Cam floor versus optimal committal regret]{Theorem~\ref{thm:t1-impquant}
in measurement: the optimal committal rule's worst-case regret coincides
with the Le~Cam floor $\tfrac{\delta}{2}(1-\TV)$ across the
total-variation sweep; at $\TV=0$ (the witness) it equals $\delta/2$, the
minimax floor of Proposition~\ref{prop:t1-minimax} with benefit gap
$\delta$. Produced by \texttt{scripts/theory\_extensions\_validation.py}
\cite{kbound2026repo}. The plot is the bridge between this section and the
next: the testing floor $1-\TV$ (this section) multiplies the benefit
margin (next section) to give worst-case regret.}
\label{fig:t1-lecam}
\end{figure}

\section{The Bayes gate and the plug-in regret identity}\label{sec:t1-gate}

The impossibility results bound what \emph{cannot} be decided; the second
result of this chapter prices decisions exactly, and in doing so converts
the abstract floor into deployed risk. The objects are the
action-conditional risks of Chapter~\ref{ch:formulation}: with
$a_{\textsc{frz}}(z)=\E[\ell(f_0(X),Y)\mid Z{=}z]$ and
$a_{\textsc{adp}}(z)=\E[\ell(f_a(X),Y)\mid Z{=}z]$, so that
$\Delta(z)=a_{\textsc{frz}}(z)-a_{\textsc{adp}}(z)$, a gate
$g\colon z\mapsto\{0,1\}$ (encoding $1=\textsc{adapt}$,
$0=\textsc{freeze}$) has risk
\[
R(g)\;=\;\E\bigl[g(Z)\,a_{\textsc{adp}}(Z)+(1-g(Z))\,a_{\textsc{frz}}(Z)\bigr].
\]
The Bayes gate is $g^\star(z)=\mathbf1[\Delta(z)>0]$; a plug-in gate built
from any label-free benefit estimator $\widehat\Delta$ is
$\widehat g(z)=\mathbf1[\widehat\Delta(z)>0]$. The paper's intuition for
what follows: a gate's excess risk over the oracle equals the benefit it
forfeits exactly where it picks the wrong sign --- near-zero-benefit
mistakes are cheap, large-benefit ones are not, which is why abstaining in
the low-margin band is safe.

\begin{theorem}[Plug-in regret decomposition]\label{thm:gate}
For every estimator $\widehat\Delta$,
\[
R(\widehat g)-R(g^\star)\;=\;\E\big[\,|\Delta(Z)|\,\mathbf1\{\sign\widehat\Delta(Z)\neq\sign\Delta(Z)\}\,\big],
\]
where $\sign(t)=\mathbf1[t>0]$. In particular $g^\star$ minimizes $R$,
regret is zero iff the gate's sign matches a.e., and on the boundary band
$\{|\Delta|<\varepsilon\}$ the committal regret is
$\le\varepsilon$---justifying \textsc{abstain} there.
\end{theorem}

\begin{proof}
We expand the paper's one-paragraph proof into its pointwise case
analysis, the integration step, and the three consequences.

\emph{Step 1 (the conditional risk is affine in the action).}
Condition on $Z=z$ and write $\rho(a;z)$ for the conditional risk of
playing action $a\in\{0,1\}$ at $z$: by definition of the
action-conditional means, $\rho(1;z)=a_{\textsc{adp}}(z)$ and
$\rho(0;z)=a_{\textsc{frz}}(z)$, so
\[
\rho(1;z)-\rho(0;z)\;=\;a_{\textsc{adp}}(z)-a_{\textsc{frz}}(z)\;=\;-\Delta(z):
\]
adapting lowers conditional risk by exactly $\Delta(z)$ (raises it when
$\Delta(z)<0$). The total risk of any gate is the average of its pointwise
choices, $R(g)=\E_Z[\rho(g(Z);Z)]$, by the tower property.

\emph{Step 2 (pointwise excess of the plug-in gate).}
Fix $z$ and compare $\widehat g(z)$ with $g^\star(z)=\mathbf1[\Delta(z)>0]$.
There are three cases.
If $\widehat g(z)=g^\star(z)$, the excess
$\rho(\widehat g(z);z)-\rho(g^\star(z);z)$ is $0$.
If $\widehat g(z)=1$ but $g^\star(z)=0$ --- the gate adapts although
$\Delta(z)\le0$ --- the excess is $\rho(1;z)-\rho(0;z)=-\Delta(z)=
|\Delta(z)|$, using $\Delta(z)\le0$.
If $\widehat g(z)=0$ but $g^\star(z)=1$ --- the gate freezes although
$\Delta(z)>0$ --- the excess is $\rho(0;z)-\rho(1;z)=\Delta(z)=
|\Delta(z)|$. In all cases,
\[
\rho(\widehat g(z);z)-\rho(g^\star(z);z)
\;=\;|\Delta(z)|\;\mathbf1\{\widehat g(z)\neq g^\star(z)\},
\]
and since $\widehat g(z)=\mathbf1[\widehat\Delta(z)>0]$ and
$g^\star(z)=\mathbf1[\Delta(z)>0]$, the disagreement event
$\{\widehat g(z)\neq g^\star(z)\}$ is exactly the sign-mismatch event
$\{\sign\widehat\Delta(z)\neq\sign\Delta(z)\}$ under the stated convention
$\sign(t)=\mathbf1[t>0]$. (Ties $\Delta(z)=0$ are immaterial: whatever the
convention assigns them, they enter the right-hand side multiplied by
$|\Delta(z)|=0$.)

\emph{Step 3 (integrate, and read off the consequences).}
Taking $\E_Z$ of the pointwise identity gives the displayed equation
exactly --- it is an identity, not an inequality, valid for every
estimator, every evidence map, and every loss with finite conditional
means; no boundedness, continuity, or distributional assumption enters.
The three consequences: (a) running the same two steps with an arbitrary
gate $g$ in place of $\widehat g$ shows
$R(g)-R(g^\star)=\E[|\Delta(Z)|\mathbf1\{g\neq g^\star\}]\ge0$, so
$g^\star$ minimizes $R$ over all gates; (b) the regret vanishes iff the
mismatch indicator vanishes $\Delta$-weighted almost everywhere, i.e.\ iff
the gate's sign matches that of $\Delta(Z)$ a.e.\ on
$\{\Delta(Z)\neq0\}$; (c) on the boundary band
$\{|\Delta(Z)|<\varepsilon\}$, the pointwise excess is
$|\Delta(z)|\,\mathbf1\{\cdot\}<\varepsilon$, so any committal choice
there --- right or wrong --- forfeits less than $\varepsilon$ of risk;
abstaining on the band is therefore safe at price at most $\varepsilon$,
which is the decision-theoretic justification for the certificate's
abstention region in Chapter~\ref{ch:theory2}.
\end{proof}

The identity composes with the matched-evidence situation of
Theorem~\ref{thm:imp} to give the floor in regret units; this is the
paper's proposition, stated verbatim.

\begin{proposition}[Label-free minimax floor]\label{prop:t1-minimax}\label{prop:gate-minimax}
In the unknowable two-point regime of Theorem~\ref{thm:imp}
($\Law(Z\mid P_T^1)=\Law(Z\mid P_T^2)$, $|\Delta^1|=|\Delta^2|=|\Delta|$,
opposite signs), every label-free estimator has worst-case expected regret
$\ge \E[|\Delta(Z)|]/2$, attained by abstaining.
\end{proposition}

\begin{proof}
Let $\widehat\Delta$ be any label-free estimator, $\widehat g$ its plug-in
gate, and allow the gate to be randomized: define
$q(z)=\Pr[\widehat g(z)=1\mid Z=z]$, the adapt-probability at evidence
$z$. Because the estimator and any internal randomization are label-free,
and the evidence law is the \emph{same} measure under both worlds, the
function $q$ is common to the two worlds --- this is Step~3 of
Theorem~\ref{thm:imp}'s proof applied conditionally at each $z$. Now
evaluate the sign-error probability at $z$ in each world: in world~1 the
benefit is negative, so the gate errs exactly when it adapts, with
probability $q(z)$; in world~2 the benefit is positive, so the gate errs
exactly when it freezes, with probability $1-q(z)$. The two error
probabilities sum to $1$ at every $z$, so the larger is $\ge\tfrac12$.
By Theorem~\ref{thm:gate}, extended to randomized gates by averaging the
pointwise identity over the gate's internal randomness, the regret in
world $j$ is
$\mathrm{reg}_j=\E\bigl[|\Delta(Z)|\cdot\Pr(\text{sign error at }Z\text{
in world }j)\bigr]$; hence
\[
\mathrm{reg}_1+\mathrm{reg}_2
=\E\bigl[|\Delta(Z)|\,(q(Z)+1-q(Z))\bigr]
=\E\bigl[|\Delta(Z)|\bigr],
\]
and therefore
$\max_j \mathrm{reg}_j\ge\tfrac12(\mathrm{reg}_1+\mathrm{reg}_2)
=\E[|\Delta(Z)|]/2$. The floor is attained: the maximally non-committal
allocation $q\equiv\tfrac12$ --- the gate-level expression of abstention,
declining to favor either sign --- has
$\mathrm{reg}_1=\mathrm{reg}_2=\E[|\Delta(Z)|]/2$, so its worst case
equals the floor exactly. In the three-action protocol the
\textsc{abstain} action realizes the same worst-case value through its
coverage accounting, which is the sense of the statement's final clause.
\end{proof}

\begin{remark}[Reading the floor]\label{rem:t1-floor}
The proposition is the regret-denominated shadow of
Theorem~\ref{thm:t1-impquant} at $\TV=0$: the testing floor says the sign
error must be $\ge\tfrac12$ in some world, and the regret identity says
each unit of sign error costs $|\Delta(Z)|$. At intermediate separations
the same composition gives worst-case regret of order
$\tfrac{|\Delta|}{2}(1-\TV)$, which is exactly the curve traced by the
validation plot of Figure~\ref{fig:t1-lecam}. Note also what the floor does
\emph{not} say: it does not make abstention free. Abstention forfeits
$\E[|\Delta|]/2$ of worst-case benefit too --- the floor is a statement
that nothing label-free can do better on matched evidence, not that
abstaining is costless. The value of abstention is realized only when the
system can \emph{also} commit on the conditions where evidence does
separate, which is the certificate's job in Chapter~\ref{ch:theory2}.
\end{remark}

\subsection{Numerical validation of the identity}\label{sec:t1-gate-val}

Because Theorem~\ref{thm:gate} is an exact identity, its validator
(\texttt{experiments/kbound/theory\_validation/val\_thm2\_regret.py}
\cite{kbound2026repo}) can demand machine precision rather than
statistical agreement, and it does, in two layers. Layer (A), the literal
content of the theorem: on $n=400{,}000$ sampled evidence points with
known conditional benefit $\Delta(z)\sim\mathcal N(0,1)$ and a non-trivial
action-conditional baseline $a_{\textsc{adp}}(z)$ (so the cancellation in
$R(\widehat g)-R(g^\star)$ is exercised, not assumed), it computes the
left side from the conditional action risks and the right side from the
mismatch identity, for plug-in estimators
$\widehat\Delta=\Delta+\mathrm{noise}$ across seven noise levels from $0$
to $2.0$; the script \emph{asserts} that the maximum absolute gap over
all noise levels is below $10^{-9}$, and the measured gap is
${\sim}10^{-17}$ --- pure floating-point rounding, as the paper reports
(``the identity holds to ${\sim}10^{-17}$''). Layer (B): the same
comparison with \emph{realized} per-sample losses drawn around the
conditional means, confirming the identity at the accuracy Monte Carlo
permits (all cells within $4$ standard errors).

The validator also checks each clause of the theorem and the proposition.
At zero estimator noise the gates agree everywhere and both sides are
exactly $0$ (regret zero iff signs match a.e.). On the boundary band
$\{|\Delta|<\varepsilon\}$ with $\varepsilon=0.05$, the maximum per-$z$
regret contribution observed is below $\varepsilon$, the ``abstaining
there costs at most $\varepsilon$'' clause. Regret rises monotonically
with estimator noise, and large noise approaches --- but does not reach
--- the coin-flip ceiling; the ceiling itself is attained by a fully
uninformative estimator (independent of $\Delta$, the empirical surrogate
of the matched-evidence regime), whose sign-mismatch rate is
${\approx}\tfrac12$ and whose regret is $0.3986$ against the floor
$\E[|\Delta|]/2=0.3985$ --- Proposition~\ref{prop:t1-minimax} realized to
three decimal places.

\begin{figure}[t]\centering
\includegraphics[width=0.62\textwidth]{fig_regret_decomposition.png}
\caption[Regret decomposition across policies]{Synthetic illustration of
the regret split (Theorem~\ref{thm:gate} /
Proposition~\ref{prop:t1-minimax}): total regret decomposes into
false-adapt, false-freeze, and abstain-coverage contributions; the
trichotomy's total is far below both trivial policies and is dominated by
benign abstain-coverage rather than harmful false-adapt. Synthetic data;
produced by \texttt{scripts/theory\_extensions\_validation.py}
\cite{kbound2026repo}. The decomposition is the deployment-facing
restatement of the regret identity: every unit of regret is attributable
to a named decision error, weighted by the benefit at stake.}
\label{fig:t1-regret}
\end{figure}

\section{Remarks: scope, mechanism, and what abstention buys}\label{sec:t1-remarks}

\begin{remark}[Relation to selective prediction]\label{rem:t1-selective}
The reject option is as old as decision theory: Chow's rule abstains from
\emph{classifying} a point when the conditional error probability exceeds
a threshold \cite{chow1970optimum}, and the selective-prediction
literature develops the resulting risk--coverage tradeoff
\cite{elyaniv2010selective, geifman2017selective}. The present chapter
lifts the same logic one level: the object declined is not a prediction
but the \emph{adaptation decision} for a deployment condition, the
``conditional error'' is the probability of picking the wrong sign of
$\Delta$, and the analogue of Chow's excess-risk computation is exactly
Theorem~\ref{thm:gate}, with $|\Delta(z)|$ playing the role of the margin.
The analogy is precise enough to transfer intuition --- low-margin points
are where rejection is cheap in both settings --- but the impossibility
half has no per-prediction counterpart: a selective classifier can always
in principle resolve its uncertainty with more labeled data, whereas
Theorem~\ref{thm:t1-impquant} exhibits decision uncertainty that no amount
of the only available (unlabeled) data resolves.
\end{remark}

\begin{remark}[The impossibility is about evidence-measurability, not
sample size]\label{rem:t1-measurability}
Every quantity in Theorem~\ref{thm:t1-impquant} is indexed by the batch
size $n$, and the witness conclusion holds \emph{for every} $n$: the floor
is $1$ at $n=10$ and at $n=10^{9}$. This locates the obstruction
precisely. Statistical folklore expects uncertainty to shrink at rate
$n^{-1/2}$, and indeed the concentration tools that power the certificate
of Chapter~\ref{ch:theory2} --- Hoeffding and empirical-Bernstein bounds
\cite{hoeffding1963probability, maurer2009empirical} --- shrink the
\emph{radius} $\varepsilon$ at exactly that rate. But a confidence radius
shrinks toward the estimand's identifiable value; when the sign of
$\Delta$ is not a function of the evidence law, there is no identifiable
value to shrink toward. Formally, the decision rule is measurable with
respect to $\sigma(Z)$ (plus independent randomness), and on the witness
pair the two worlds induce the \emph{same} measure on $\sigma(Z)$; the
answer is not a random variable on the space the rule lives on. More data
refines the empirical estimate of that common measure --- it cannot break
the tie. This is why the chapter's results are phrased as
non-identifiability rather than as a rate: the failure is of the
$\sigma$-algebra, not of the variance.
\end{remark}

\begin{remark}[What breaks without each hypothesis]\label{rem:t1-necessity}
It is worth walking the hypotheses of the two theorems once more, now with
their proofs in hand. \emph{Matched evidence laws} (hypothesis (i)): drop
it and Lemma~\ref{lem:t1-affinity} immediately monetizes the difference
--- the floor falls to $1-\TV<1$, and a likelihood-ratio test on the
evidence achieves it; this is not a pathology but the design target of
evidence engineering (Remark~\ref{rem:t1-tightness}). \emph{Opposite
benefit signs} (hypothesis (ii)): drop it and the matched evidence is
harmless, since one committed action is correct in both worlds; the
trichotomy collapses to a dichotomy on such families. \emph{Fixed, known
$(f_0,f_a)$} (Assumption~\ref{asm:form-triple}): if the candidate were not
fixed at decision time, $Z$ would not be a world-independent function of
the covariates and Step~2 of Theorem~\ref{thm:imp}'s proof would fail to
produce matched evidence laws --- but the conclusion would not improve,
because the benefit being certified would itself be a moving target.
\emph{The label-free protocol} (Assumption~\ref{asm:form-protocol}): a
single labeled target batch breaks the construction completely --- the
per-sample paired benefits $\delta_i$ become observable, their mean
estimates $\Delta$ directly, and the witness worlds separate at the first
label. The impossibility is exactly as strong as the constraint that
creates it, which is why the paper is careful to state the protocol
(Remark~\ref{rem:form-labelfree}) before the theorems.
\end{remark}

\begin{remark}[What abstention buys, and what comes next]\label{rem:t1-buys}
The two results of this chapter are deliberately one-sided: they establish
that on matched-evidence pairs every committal policy pays the floor, and
that the unique worst-case-optimal action there is \textsc{abstain}. What
they do not yet provide is a \emph{computable} rule that abstains there
and commits elsewhere with guarantees. That is the certificate's job:
Theorem~\ref{thm:cert} (Chapter~\ref{ch:theory2}) shows that an interval
$\widehat\Delta\pm\varepsilon$ with $1-\alpha$ coverage yields an
adapt/freeze/abstain rule whose false-adapt and false-freeze rates are
each at most $\alpha$; its forced-abstention corollary closes the loop
with this chapter by showing that on matched evidence any such valid rule
\emph{must} abstain with probability at least $1-2\alpha$ --- abstention
is mandatory, not merely prudent; and its sequential complement
(Theorem~\ref{thm:evalue}) extends the guarantee to streams with optional
stopping. On the positive side of the boundary,
Theorem~\ref{thm:pos} proves that under covariate shift with a bounded
density ratio the sign of $\Delta$ \emph{is} identifiable from labeled
source plus unlabeled target \cite{shimodaira2000improving}, and
Theorem~\ref{thm:disagree} characterizes $\sign\Delta$ in general as an
ordinal accuracy comparison on the observable disagreement region $D$ ---
the quantity whose label-free bracketing is the program's one remaining
open piece (Conjecture~\ref{conj:gen}). The unknowable regime proved here
is thus not the end of the story but its boundary condition: everything
Chapter~\ref{ch:theory2} certifies, it certifies \emph{because} it
abstains where this chapter says it must.
\end{remark}
